%==============================================================================
% INTRODUCTION GÉNÉRALE
%==============================================================================
\chapter*{Introduction Générale}
\addcontentsline{toc}{chapter}{Introduction Générale}

La Business Intelligence (BI) constitue aujourd'hui un pilier fondamental de la prise de décision dans les organisations de toute taille. Elle englobe l'ensemble des processus de collecte, de transformation et d'analyse des données visant à produire des informations actionnables pour les décideurs. Toutefois, les outils traditionnels de BI --- tels que Power BI, Tableau ou Qlik --- imposent une courbe d'apprentissage significative, nécessitant des compétences en modélisation des données, en langages de requête (SQL, DAX, M) et en conception de visualisations.

Parallèlement, l'émergence des modèles de langage à grande échelle (LLM) et des systèmes multi-agents ouvre des perspectives radicalement nouvelles pour démocratiser l'accès à l'analyse de données. Ces modèles, capables de comprendre le langage naturel et de raisonner sur des schémas de bases de données, permettent d'envisager un nouveau paradigme d'interaction avec les systèmes décisionnels : le \textit{Generative BI}, où l'utilisateur exprime simplement son besoin en langage naturel, et le système se charge automatiquement de la formulation des requêtes, de l'agrégation des données et de la mise en forme des résultats.

C'est dans ce contexte que s'inscrit notre projet de fin d'année : la conception et la réalisation de \textbf{SmartBI}, une plateforme de Business Intelligence pilotée par des agents d'intelligence artificielle. L'ambition est de créer un système capable de recevoir un fichier de données brut (CSV ou XLSX), de l'analyser automatiquement, de proposer des indicateurs clés de performance pertinents, de générer les requêtes SQL nécessaires et de produire un tableau de bord interactif complet --- le tout sans intervention humaine experte. L'utilisateur peut également dialoguer avec le système en langage naturel pour obtenir des analyses ciblées.

Pour atteindre cet objectif, nous avons adopté une architecture multi-agents fondée sur le framework \textbf{LangGraph}, où trois agents spécialisés --- un analyste métier, un ingénieur de données et un designer --- collaborent de manière séquentielle avec un mécanisme de rétroaction et de correction automatique d'erreurs. Les modèles de langage sont exécutés localement via \textbf{Ollama}, ce qui garantit la souveraineté des données et une latence acceptable.

Le présent rapport est structuré en trois chapitres. Le \textbf{premier chapitre} pose le contexte général du projet, expose la problématique à laquelle il répond, dresse un état de l'art des solutions existantes et présente les objectifs visés. Le \textbf{deuxième chapitre} détaille l'analyse fonctionnelle et la conception technique du système, incluant l'architecture globale, le pipeline d'agents, la modélisation des données et la conception de l'interface utilisateur. Le \textbf{troisième chapitre} est consacré à la réalisation : il décrit l'environnement de développement, présente les interfaces réalisées, illustre le fonctionnement du système à travers des scénarios d'utilisation et discute les difficultés rencontrées. Nous concluons par une synthèse des travaux réalisés et une discussion sur les perspectives d'évolution.
